\documentclass[11pt]{article}

\usepackage{amsmath,amssymb,amsthm}
\usepackage{xcolor}
\usepackage{xr}
\externaldocument{main-paper-labels}

\title{Online Supplement to ``On the Network Coding Capacity of Wireless Multicast Networks under One-Hop Interference''}
\author{Xiaohong Cai and Raymond W. Yeung}
\date{}

\newtheorem{lemma}{Lemma}
\newtheorem{claim}{Claim}[lemma]
\newtheorem{remark}{Remark}
\newtheorem{theorem}{Theorem}
\newtheorem{definition}{Definition}
\newtheorem{corollary}{Corollary}


\begin{document}

\maketitle

\section{Finite Matching-Based Verification}

For each reduced butterfly-like configuration, we construct the associated directed hypergraph \(H=(\mathcal N,\mathcal E_H)\). A feasible simultaneous transmission set is a matching of \(H\), i.e., a set of hyperedges no two of which are incident to a common node. Let \(\mathcal F_H\) denote the set of all matchings of \(H\), including the empty matching, and let \(\omega_\sigma\) be the fraction of the resource assigned to \(\sigma\in\mathcal F_H\).

For each destination \(t\in\{X,Y\}\), let \(\mathcal C_H^{st}\) denote the set of all \(s\)-\(t\) cutsets in \(H\), as defined in the main paper. Thus, for an \(s\)-\(t\) cut with source side \(\mathcal S\subseteq\mathcal N\), where \(s\in\mathcal S\) and \(t\in\mathcal N\setminus\mathcal S\), the associated cutset contains every hyperedge \((u,V)\in\mathcal E_H\) such that \(u\in\mathcal S\) and \(V\cap(\mathcal N\setminus\mathcal S)\neq\emptyset\).

Following the notation for butterfly-like networks in the main paper, the subpaths \(P^{AI}\), \(P^{BI}\), \(P^{DX}\), \(P^{DY}\), and \(P^{ID}\) are denoted by \(c\), \(d\), \(g\), \(h\), and \(i\), respectively. We define \(\mathcal E_{\mathrm{inner}}\) as the set of all directed edges on \(c\), \(d\), \(g\), and \(h\), together with those on \(i\) when \(I\neq D\).

For each reduced configuration, we maximize the multicast rate \(R\) over the matching weights \(\{\omega_\sigma\}_{\sigma\in\mathcal F_H}\). Let \(C_{(u,V)}\) denote the capacity of hyperedge \((u,V)\) induced by these weights. We solve the following linear program:
\[
\begin{aligned}
\text{maximize}\quad
& R\\
\text{subject to}\quad
& \sum_{\sigma\in\mathcal F_H}\omega_\sigma=1,\\
& C_{(u,V)}=\sum_{\sigma\in\mathcal F_H:(u,V)\in\sigma}\omega_\sigma,
&& (u,V)\in\mathcal E_H,\\
& R\le\sum_{(u,V)\in K}C_{(u,V)},
&& \substack{K\in\mathcal C_H^{st}, \;\; t\in\{X,Y\}},\\
& \sum_{\substack{\sigma\in\mathcal F_H:\,\exists (u,V)\in\sigma,\\ v\in V}}\omega_\sigma=\frac13,
&& (u,v)\in\mathcal E_{\mathrm{inner}},\\
& \omega_\sigma\ge0,
&& \sigma\in\mathcal F_H,\\
& 0\le R\le1.
\end{aligned}
\]
The equality constraint indexed by \((u,v)\in\mathcal E_{\mathrm{inner}}\) fixes the induced capacity of every directed edge on an inner subpath to \(1/3\).

The exhaustive computation verifies every reduced configuration for both \(I=D\) and \(I\neq D\). In each case, the configurations with LP value at least \(5/6\) are exactly those satisfying the corresponding structural conditions in the main paper. Together with the \(5/6\) upper bound for butterfly-like networks, this shows that every reduced configuration satisfying these conditions has capacity exactly \(5/6\). The implementation and complete machine-readable output are provided with this supplement.

\section{Proof of Theorem~\ref{theo: 2-and-3 paths}}\label{Appendix: Proof_of_three_path_disjoint}

Let \(\mathcal T=\{X,Y\}\). Without loss of generality, assume that \(G_0\) contains two disjoint routing paths \(P_X\) and \(\bar P_X\) from \(s\) to \(X\), and three disjoint routing paths \(Q_1\), \(Q_2\), and \(Q_3\) from \(s\) to \(Y\).
We use the subpath notation \(P^{uv}\) defined in Section~\ref{sec: Structural Preliminaries} in the main paper. For ease of notation, when \(u=v\), \(P^{uv}\) is omitted in path constructions.

Since \(G_0\) is a directed acyclic graph, its nodes admit a topological ordering. Fix one such ordering and denote it by \(``<"\). For any two distinct nodes \(u,v\in N(G_0)\), if there is a directed path from \(u\) to \(v\), then \(u<v\). Throughout the proof, this ordering is used both to compare nodes and to identify the last node in a node set.

The proof below constructs four routing paths: two disjoint paths connecting \(s\) to \(X\), and two disjoint paths connecting \(s\) to \(Y\). For each construction, let \(G_r\) be the subgraph whose node set is the union of the node sets of the four selected paths and whose edge set is the union of their edge sets. By Definition~\ref{def: routing graph}, \(G_r\) is a routing graph. Since all selected paths lie in \(G_0\), \(G_r\) is a routing subgraph of \(G_0\). It remains to show that this routing subgraph contains no intersection node. Since the two paths to the same destination are disjoint by construction, we only need to check those pairs consisting of one path to \(X\) and one path to \(Y\). By Propositions~\ref{prop: intersection if not disjoint and multicast} and~\ref{prop: Multicast have no intersection node}, it is enough to verify that each such pair is either disjoint or a multicast.

For each \(i\in\{1,2,3\}\), let \(M_i\) be the last node in
\begin{equation}\label{eq: 1}
    N(Q_i)\cap\bigl(N(P_X)\cup N(\bar P_X)\bigr).
\end{equation}
If \(M_i\neq s\), then \(M_i\) lies on exactly one of \(P_X\) and \(\bar P_X\), since \(P_X\) and \(\bar P_X\) are disjoint except at \(s\) and \(X\), and a routing path to \(Y\) cannot pass through the destination \(X\).

We distinguish three cases according to how many of \(M_1\), \(M_2\), and \(M_3\) are equal to \(s\) and whether the remaining nodes lie on the same path in \(\{P_X,\bar P_X\}\).

\smallskip
\noindent\textbf{Case 1.}
At least two of \(M_1\), \(M_2\), and \(M_3\) are equal to \(s\). Without loss of generality, let \(M_1=M_2=s\). Then neither \(Q_1\) nor \(Q_2\) has a node in \(N(P_X)\cup N(\bar P_X)\) other than \(s\). Thus
\[
P_X, \bar P_X, Q_1, Q_2
\]
already form the required four paths.
Every pair connecting to different destinations is disjoint, so the subgraph formed by these four paths contains no intersection node.

\smallskip
\noindent\textbf{Case 2.}
Case 1 does not occur, and there exist two non-source nodes among \(M_1\), \(M_2\), and \(M_3\) that lie on different paths in \(\{P_X,\bar P_X\}\). That is, there exist \(i\neq j\) such that \(M_i,M_j\neq s\), with \(M_i\) and \(M_j\) lying on different paths in \(\{P_X,\bar P_X\}\). Without loss of generality, assume that \(M_i\in N(P_X)\) and \(M_j\in N(\bar P_X)\). Define
\[
Q_i'=(P_X^{sM_i},Q_i^{M_iY}),
\qquad
Q_j'=(\bar P_X^{sM_j},Q_j^{M_jY}).
\]
\par\vspace{0.4em}
\noindent\emph{Disjointness of \(Q_i'\) and \(Q_j'\)}:
The subpaths \(P_X^{sM_i}\) and \(\bar P_X^{sM_j}\) are disjoint because \(P_X\) and \(\bar P_X\) are disjoint. Likewise, \(Q_i^{M_iY}\) and \(Q_j^{M_jY}\) are disjoint because \(Q_i\) and \(Q_j\) are disjoint. It remains to consider the two cross-pairs. By the definition of \(M_i\), \(Q_i^{M_iY}\) shares no node with \(\bar P_X^{sM_j}\); similarly, by the definition of \(M_j\), \(Q_j^{M_jY}\) shares no node with \(P_X^{sM_i}\). Thus, \(Q_i'\) and \(Q_j'\) are disjoint.

\par\vspace{0.4em}
\noindent\emph{Multicast pairs formed by \(P_X\) and \(Q_i'\), and by \(\bar P_X\) and \(Q_j'\)}:
The paths \(P_X\) and \(Q_i'\) share the initial subpath \(P_X^{sM_i}\), and \(Q_i^{M_iY}\) shares no further node with \(P_X\) by the definition of \(M_i\). Thus, \(P_X\) and \(Q_i'\) form a multicast. Similarly, \(\bar P_X\) and \(Q_j'\) form a multicast with common prefix \(\bar P_X^{sM_j}\).

\par\vspace{0.4em}
\noindent\emph{Disjointness of \(P_X\) from \(Q_j'\) and of \(\bar P_X\) from \(Q_i'\)}:
The subpath \(\bar P_X^{sM_j}\) is disjoint from \(P_X\), and \(Q_j^{M_jY}\) shares no node with \(P_X\) by the definition of \(M_j\). Similarly, the subpath \(P_X^{sM_i}\) is disjoint from \(\bar P_X\), and \(Q_i^{M_iY}\) shares no node with \(\bar P_X\) by the definition of \(M_i\).

\smallskip
\noindent Consequently, \(P_X\), \(\bar P_X\), \(Q_i'\), and \(Q_j'\) form a routing subgraph of \(G_0\), and every pair of paths connecting to different destinations is either disjoint or a multicast. By Propositions~\ref{prop: intersection if not disjoint and multicast} and~\ref{prop: Multicast have no intersection node}, this routing subgraph contains no intersection node.

\smallskip
\noindent\textbf{Case 3.}
Neither Case 1 nor Case 2 occurs. Since Case 1 does not occur, at most one of \(M_1\), \(M_2\), and \(M_3\) is \(s\). Since Case 2 does not occur, all \(M_i\) that are non-source nodes lie on the same path in \(\{P_X,\bar P_X\}\). Moreover, these non-source nodes are distinct because \(Q_1\), \(Q_2\), and \(Q_3\) are disjoint (cf.~\eqref{eq: 1}). Without loss of generality, assume that this same path is \(P_X\), and relabel \(Q_1,Q_2\), and \(Q_3\) so that
\[
M_1<M_2<M_3
\qquad\text{on }P_X,
\]
where \(M_1=s\) is allowed.

For each \(i\), let \(L_i\) be the last node in
\begin{equation}\label{eq: 2}
    N(Q_i)\cap N(\bar P_X).
\end{equation}
If \(Q_i\) and \(\bar P_X\) share no non-source node, then \(L_i=s\).

\par\vspace{0.4em}
\noindent\textbf{\emph{Subcase 3a.}}
There exist \(i<j\) such that \(L_j\le L_i\). Define
\[
P_X'=(Q_j^{sM_j},P_X^{M_jX}),
\qquad
\bar P_X'=(Q_i^{sL_i},\bar P_X^{L_iX}),
\]
\par\vspace{0.4em}
\noindent\emph{Disjointness of \(P_X'\) and \(\bar P_X'\)}:
The subpaths \(Q_j^{sM_j}\) and \(Q_i^{sL_i}\) are disjoint because \(Q_j\) and \(Q_i\) are disjoint. Likewise, \(P_X^{M_jX}\) and \(\bar P_X^{L_iX}\) are disjoint because \(P_X\) and \(\bar P_X\) are disjoint. It remains to consider the two cross-pairs. Since \(M_i<M_j\) and \(M_i\) is the last node in \(N(Q_i)\cap N(P_X)\), we have \(N(Q_i)\cap N(P_X^{M_jX})=\nobreak\emptyset\). Therefore, \(Q_i^{sL_i}\) shares no node with \(P_X^{M_jX}\). Since \(Q_i\) and \(Q_j\) are disjoint, \(L_i=L_j\) can occur only when \(L_i=L_j=s\). Together with \(L_j\le L_i\) and the fact that \(L_j\) is the last node in \(N(Q_j)\cap N(\bar{P}_X)\), this gives
\[
\bigl(N(Q_j)\setminus\{s\}\bigr)\cap N(\bar P_X^{L_iX})=\emptyset.
\]
Therefore, \(Q_j^{sM_j}\) shares no node with \(\bar P_X^{L_iX}\) other than possibly \(s\). Thus, \(P_X'\) and \(\bar P_X'\) are disjoint.

\par\vspace{0.4em}
\noindent\emph{Disjointness of \(\bar P_X'\) from \(Q_j\) and of \(P_X'\) from \(Q_i\)}:
\(Q_j\) is disjoint from \(Q_i^{sL_i}\), and \(\bigl(N(Q_j)\setminus\{s\}\bigr)\cap N(\bar P_X^{L_iX})=\nobreak\emptyset\), as shown above. Moreover, \(Q_i\) is disjoint from \(Q_j^{sM_j}\), and \(N(Q_i)\cap N(P_X^{M_jX})=\nobreak\emptyset\), as shown above.

\par\vspace{0.4em}
\noindent\emph{Multicast pair formed by \(P_X'\) and \(Q_j\)}:
The paths \(P_X'\) and \(Q_j\) share the initial subpath \(Q_j^{sM_j}\), and \(Q_j^{M_jY}\) shares no further node with \(P_X\) by the definition of \(M_j\) (cf.~\eqref{eq: 1}). Thus, \(P_X'\) and \(Q_j\) form a multicast.

\par\vspace{0.4em}
\noindent\emph{Analysis of \(\bar P_X'\) and \(Q_i\)}:
If \(L_i\neq s\), then \(\bar P_X'\) and \(Q_i\) share the initial subpath \(Q_i^{sL_i}\), and the remainder of \(Q_i\) shares no further node with \(\bar P_X\) by the definition of \(L_i\) (cf.~\eqref{eq: 2}); hence, they form a multicast. If \(L_i=s\), then \(\bar P_X'\) and \(Q_i\) are disjoint.
\par\vspace{0.6em}
\noindent Since \(Q_i\) and \(Q_j\) are disjoint, \(P_X'\), \(\bar P_X'\), \(Q_i\), and \(Q_j\) form a routing subgraph of \(G_0\), and every pair of paths connecting to different destinations is either disjoint or a multicast. By Propositions~\ref{prop: intersection if not disjoint and multicast} and~\ref{prop: Multicast have no intersection node}, this routing subgraph contains no intersection node.

\par\vspace{0.4em}
\noindent\textbf{\emph{Subcase 3b.}}
No pair \(i<j\) satisfies \(L_j\le L_i\). This implies
\[
L_1<L_2<L_3
\qquad\text{on }\bar P_X.
\]

Consider \(Q_3^{sM_3}\), where \(M_3\in N(P_X)\), and let \(A_0=M_3\). By tracing \(Q_3\) backwards from \(A_0\), we will construct nodes
\[
A_0,A_1,A_2,\ldots,A_l\in N(P_X)\cap N(Q_3^{sM_3})\]
and
\[
B_1,B_2,\ldots,B_l\in N(\bar P_X)\cap N(Q_3^{sM_3}),
\]
where \(l\ge 1\), such that
\[
M_3=A_0>B_1>A_1>B_2>\cdots > A_{l-1}>B_l\ge A_l =s.
\]
Specifically, given \(A_{k-1}\neq s\), let \(B_k\) be the last node in \(N(Q_3^{sA_{k-1}})\cap N(\bar P_X)\).\footnote{If \(Q_3^{sA_{k-1}}\) and \(\bar P_X\) share no non-source node, then \(B_k=s\).}  Then let \(A_k\) be the last node in \(N(Q_3^{sB_k})\cap N(P_X)\).\footnote{If \(Q_3^{sB_k}\) and \(P_X\) share no non-source node, then \(A_k=s\).}  By construction, for each \(k\) with \(A_{k-1}\neq s\), we have \(A_{k-1}>B_k\ge A_k\), where equality can occur only when \(B_k=A_k=s\). Consequently, \(A_0>A_1>A_2>\cdots > A_l=s\) on \(P_X\).
For this sequence, since
\begin{enumerate}
    \item \(A_0=M_3>M_2\);
    \item the sequence eventually reaches \(s\);
    \item no \(A_k\) can equal \(M_2\) since \(Q_2\) and \(Q_3\) are disjoint (note that \(A_k \in N(Q_3)\) and \(M_2\in N(Q_2)\)),
\end{enumerate}
there exists an \(r\ge 1\) such that
\[
A_r<M_2<A_{r-1}
\qquad\text{on }P_X.
\]
Since \(L_1<L_2\), we have \(L_2\neq s\). Moreover, \(B_r\neq L_2\) because equality would make \(Q_2\) and \(Q_3\) share the non-source node \(L_2\) (note that \(B_r\in N(Q_3)\) and \(L_2\in N(Q_2)\)). Therefore, either \(B_r>L_2\) or \(B_r<L_2\) on \(\bar P_X\).

\par\vspace{0.4em}
\noindent\textbf{\emph{Subcase 3b(i).}}
\(B_r>L_2\). Define
\[
P_X'=(Q_2^{sM_2},P_X^{M_2X}),
\qquad
\bar P_X'=(Q_3^{sB_r},\bar P_X^{B_rX}),
\]
\par\vspace{0.4em}
\noindent\emph{Disjointness of \(P_X'\) and \(\bar P_X'\)}:
The subpaths \(Q_2^{sM_2}\) and \(Q_3^{sB_r}\) are disjoint because \(Q_2\) and \(Q_3\) are disjoint. Likewise, \(P_X^{M_2X}\) and \(\bar P_X^{B_rX}\) are disjoint because \(P_X\) and \(\bar P_X\) are disjoint. It remains to consider the two cross-pairs. Since \(L_2<B_r\) on \(\bar{P}_X\) and \(L_2\) is the last node in \(N(Q_2)\cap N(\bar P_X)\), we have \(N(Q_2)\cap N(\bar P_X^{B_rX})=\nobreak\emptyset\). Therefore, \(Q_2^{sM_2}\) shares no node with \(\bar P_X^{B_rX}\). Since \(A_r<M_2\) and \(A_r\) is the last node in \(N(Q_3^{sB_r})\cap N(P_X)\), we have \(N(Q_3^{sB_r})\cap N(P_X^{M_2X})=\nobreak\emptyset\). Therefore, \(Q_3^{sB_r}\) shares no node with \(P_X^{M_2X}\). Thus, \(P_X'\) and \(\bar P_X'\) are disjoint.

\par\vspace{0.4em}
\noindent\emph{Disjointness of \(\bar P_X'\) from \(Q_2\)}:
\(Q_2\) is disjoint from \(Q_3^{sB_r}\). Moreover, \(N(Q_2)\cap N(\bar P_X^{B_rX})=\nobreak\emptyset\) as shown above.

\par\vspace{0.4em}
\noindent\emph{Disjointness of \(P_X'\) from \(Q_1\)}:
\(Q_1\) is disjoint from \(Q_2^{sM_2}\). Moreover, \(N(Q_1)\cap N(P_X^{M_2X})=\nobreak\emptyset\), since \(M_1<M_2\) and \(M_1\) is the last node in \(N(Q_1)\cap N(P_X)\).

\par\vspace{0.4em}
\noindent\emph{Disjointness of \(\bar P_X'\) from \(Q_1\)}:
\(Q_1\) is disjoint from \(Q_3^{sB_r}\). Moreover, \(N(Q_1)\cap N(\bar P_X^{B_rX})=\nobreak\emptyset\), since \(L_1<L_2<B_r\) and \(L_1\) is the last node in \(N(Q_1)\cap N(\bar P_X)\).

\par\vspace{0.4em}
\noindent\emph{Multicast pair formed by \(P_X'\) and \(Q_2\)}:
The paths \(P_X'\) and \(Q_2\) share the initial subpath \(Q_2^{sM_2}\), and \(Q_2^{M_2Y}\) shares no further node with \(P_X\) by the definition of \(M_2\) (cf.~\eqref{eq: 1}). Thus, \(P_X'\) and \(Q_2\) form a multicast.

\par\vspace{0.6em}
\noindent Since \(Q_2\) and \(Q_1\) are disjoint, \(P_X'\), \(\bar P_X'\), \(Q_2\), and \(Q_1\) form a routing subgraph of \(G_0\), and every pair of paths connecting to different destinations is either disjoint or a multicast. By Propositions~\ref{prop: intersection if not disjoint and multicast} and~\ref{prop: Multicast have no intersection node}, this routing subgraph contains no intersection node.

\par\vspace{0.4em}
\noindent\textbf{\emph{Subcase 3b(ii).}}
\(B_r<L_2\). Define
\[
P_X'=(Q_3^{sA_{r-1}},P_X^{A_{r-1}X}),
\qquad
\bar P_X'=(Q_2^{sL_2},\bar P_X^{L_2X}),
\]
\par\vspace{0.4em}
\noindent\emph{Disjointness of \(P_X'\) and \(\bar P_X'\)}:
The subpaths \(Q_3^{sA_{r-1}}\) and \(Q_2^{sL_2}\) are disjoint because \(Q_3\) and \(Q_2\) are disjoint. Likewise, \(P_X^{A_{r-1}X}\) and \(\bar P_X^{L_2X}\) are disjoint because \(P_X\) and \(\bar P_X\) are disjoint. It remains to consider the two cross-pairs. Since \(B_r<L_2\) on \(\bar{P}_X\) and \(B_r\) is the last node in \(N(Q_3^{sA_{r-1}})\cap N(\bar P_X)\), we have \(N(Q_3^{sA_{r-1}})\cap N(\bar P_X^{L_2X})=\nobreak\emptyset\). Therefore, \(Q_3^{sA_{r-1}}\) shares no node with \(\bar P_X^{L_2X}\). Since \(M_2<A_{r-1}\) and \(M_2\) is the last node in \(N(Q_2)\cap N(P_X)\), we have \(N(Q_2)\cap N(P_X^{A_{r-1}X})=\nobreak\emptyset\). Therefore, \(Q_2^{sL_2}\) shares no node with \(P_X^{A_{r-1}X}\). Thus, \(P_X'\) and \(\bar P_X'\) are disjoint.

\par\vspace{0.4em}
\noindent\emph{Disjointness of \(P_X'\) from \(Q_2\)}:
\(Q_2\) is disjoint from \(Q_3^{sA_{r-1}}\), and \(N(Q_2)\cap N(P_X^{A_{r-1}X})=\nobreak\emptyset\), as shown above.

\par\vspace{0.4em}
\noindent\emph{Disjointness of \(P_X'\) from \(Q_1\)}:
\(Q_1\) is disjoint from \(Q_3^{sA_{r-1}}\), and \(N(Q_1)\cap N(P_X^{A_{r-1}X})=\nobreak\emptyset\), since \(M_1<M_2<A_{r-1}\) and \(M_1\) is the last node in \(N(Q_1)\cap N(P_X)\).

\par\vspace{0.4em}
\noindent\emph{Disjointness of \(\bar P_X'\) from \(Q_1\)}:
\(Q_1\) is disjoint from \(Q_2^{sL_2}\). Moreover, \(N(Q_1)\cap N(\bar P_X^{L_2X})=\nobreak\emptyset\), since \(L_1<L_2\) and \(L_1\) is the last node in \(N(Q_1)\cap N(\bar P_X)\).

\par\vspace{0.4em}
\noindent\emph{Multicast pair formed by \(\bar P_X'\) and \(Q_2\)}:
The paths \(\bar P_X'\) and \(Q_2\) share the initial subpath \(Q_2^{sL_2}\), and the remainder of \(Q_2\) shares no further node with \(\bar P_X\) by the definition of \(L_2\) (cf.~\eqref{eq: 2}). Thus, \(\bar P_X'\) and \(Q_2\) form a multicast.

\par\vspace{0.6em}
\noindent Since \(Q_2\) and \(Q_1\) are disjoint, \(P_X'\), \(\bar P_X'\), \(Q_2\), and \(Q_1\) form a routing subgraph of \(G_0\), and every pair of paths connecting to different destinations is either disjoint or a multicast. By Propositions~\ref{prop: intersection if not disjoint and multicast} and~\ref{prop: Multicast have no intersection node}, this routing subgraph contains no intersection node.

\par\vspace{0.6em}
In all cases, we obtain two disjoint routing paths in \(G_0\) to each destination, and the subgraph formed by these paths contains no intersection node. This completes the proof.

\end{document}
